\chapter{认知运行时}

\section{运行时状态}

认知运行时组织当前可供推理和调用方使用的认知状态。运行时可以丢失并从权威数据重建，任何持久认知真值都不依赖进程内缓存存活。

\begin{description}
\item[会话（Session）] 一段连续运行范围，不等同于聊天线程或经历片段。
\item[活动认知（Active Cognition）] 当前较容易重新激活的认知引用和少量运行元数据。
\item[工作上下文（WorkContext）] 可暂停、恢复和结束的任务或问题上下文；同一主体可以同时存在多个工作上下文。
\item[驻留集（ResidentSet）] 当前运行时直接保留的认知引用集合。淘汰只改变运行时可用性，不修改持久认知。
\item[调用方工作集（ConsumerWorkingSet）] 按调用方能力、意图和预算选择的有界认知集合。
\item[受管理上下文（Managed Context）] 为模型或其他有状态消费者维护的上下文同步状态。
\end{description}

\section{注意与即时评价}

注意分配在有限预算下选择当前需要处理的信息、工作上下文和认知对象。即时评价根据自我、社会认知、动机和当前环境估计相关性、机会、风险、可控性等运行时性质。

注意和即时评价主要属于运行时。只有形成稳定、可追溯且需要跨会话长期保留的倾向时，才通过相应领域提案进入自我、社会认知、动机或记忆。

\section{使用事件}

长期使用事件记录调用方对某项认知的真实使用及后续结果。候选生成、检索命中、排序和过滤属于检索轨迹，不直接成为长期使用事件。

使用事件可以区分呈现、引用、据此行动、结果支持、结果反驳、纠正和显式固定等语义。仅进入调用方上下文的呈现不等于有意义使用，也不自动改变长期关联或可访问性。

每个需要持久记录的使用事件由调用方提供稳定事件标识。重试同一事件必须得到幂等结果；相同事件标识携带不同内容时返回冲突。使用事件与由其直接触发的运行时更新保持一致提交。

\section{可访问性}

可访问性描述在当前线索和上下文中取得某项认知的难易程度。自动可访问性是查询时派生状态，可以考虑认知年龄、类型化使用、当前线索、工作上下文和查询预算。

自动可访问性不修改认知真值和生命周期。显式固定、精确访问限制、抑制和清除分别由独立策略或管理状态表达。自动策略不得仅依据固定年龄把仍然合法存在的认知永久变成不可恢复对象。

\section{受管理上下文}

对有状态消费者追加上下文前，系统核对调用方身份、策略、提示前缀、来源修订、访问控制、对象生命周期和渲染合同。上述条件发生不兼容变化时重建上下文；满足兼容条件时可以追加。上下文缓存永远不代替认知权威。
